\documentclass[11pt,letterpaper]{article}

% ---------- Packages (all ship with Overleaf's default TeX Live) ----------
\usepackage[margin=1in]{geometry}
\usepackage{amsmath}
\usepackage{graphicx}
\usepackage{booktabs}
\usepackage{pgfplotstable}
\usepackage{float}
\usepackage[font=small,labelfont=bf]{caption}
\usepackage{setspace}
\usepackage[colorlinks=true,allcolors=black]{hyperref}

\pgfplotsset{compat=1.16}

% ---------- Dependencies: edit these two lines if your filenames differ ----------
\newcommand{\datafile}{results.csv}
\newcommand{\figfile}{scaling.png}

\setlength{\parindent}{1.5em}
\setlength{\parskip}{0pt}

\begin{document}

% ---------- Header block ----------
\noindent
Author: Trevor Robbins\\
Class: CS431\\
Professor: Choi\\
Date: 10/30/26

\begin{center}
  {\large\bfseries Analysis of Prefix-Sum Serial vs.\ Parallel Algorithms\\
  (Hillis--Steele and Blelloch Scan)}
\end{center}

\onehalfspacing

% ---------- Introduction ----------
In this homework assignment, I demonstrate the performance differences between
three implementations of the prefix sum. The prefix sum is an operator that,
given an array of size $n$, iterates up to $n$ while calculating a running sum
across the elements. For this case study, I use a serial version as the control
and compare it against the Hillis--Steele and Blelloch algorithms, with the
number of CPU threads as an adjustable parameter set to 2, 4, 8, 16, 32, and 64.

\section*{Serial vs.\ Hillis--Steele vs.\ Blelloch Scan}

The serial implementation is practically identical to the definition above: the
iterator runs from $1$ to $n$, carrying the running total forward until the
element at $n$ has been summed. This serial algorithm has bottlenecks that
parallelization can capitalize on, which is where the Hillis--Steele parallel
algorithm comes in. Hillis--Steele iterates over the array while doubling the
stride (offset) at each step, so each element's partial sum propagates
logarithmically across the array. In turn, this algorithm performs an extra
$\log_2(n)$ factor of work per element. With $n$ processors this amortizes to
$O(\log n)$ time; however, because we only use up to 64 CPU threads, the cost
simplifies to $O(n \log n)$ total work.

Finally, there is the Blelloch scan. This algorithm builds the running sum with
an up-sweep (a tree of partial sums) followed by a down-sweep (propagating those
sums back down the tree). The tree has a depth of $\log_2(n)$, and the total
work is $2(n-1)$ additions, giving $O(n)$ complexity.

% ---------- Analysis ----------
\section*{Analysis}

Table~\ref{tab:results} lists the mean runtimes over 20 runs and the resulting
speedups for each thread count, and Figure~\ref{fig:results} plots them.

\begin{table}[H]
  \centering
  \caption{Mean runtime (seconds) and speedup relative to the serial baseline,
  averaged over 20 runs per thread count.}
  \label{tab:results}
  \resizebox{\textwidth}{!}{%
  \pgfplotstabletypeset[
    col sep=comma,
    every head row/.style={before row=\toprule, after row=\midrule},
    every last row/.style={after row=\bottomrule},
    columns/threads/.style={column name={Threads}, int detect, column type=c},
    columns/runs/.style={column name={Runs}, int detect, column type=c},
    columns/serial_s/.style={column name={Serial (s)}, fixed, fixed zerofill, precision=4, column type=c},
    columns/hillis_steele_s/.style={column name={Hillis--Steele (s)}, fixed, fixed zerofill, precision=4, column type=c},
    columns/blelloch_s/.style={column name={Blelloch (s)}, fixed, fixed zerofill, precision=4, column type=c},
    columns/hillis_steele_speedup/.style={column name={Hillis--Steele Speedup}, fixed, fixed zerofill, precision=3, column type=c},
    columns/blelloch_speedup/.style={column name={Blelloch Speedup}, fixed, fixed zerofill, precision=3, column type=c},
  ]{\datafile}%
  }
\end{table}

\begin{figure}[H]
  \centering
  \includegraphics[width=\textwidth]{\figfile}
  \caption{Left: mean prefix-sum runtime versus thread count (log scale).
  Right: speedup relative to the serial implementation (log scale; 1 = serial).}
  \label{fig:results}
\end{figure}

Reading the results above, the serial implementation stays controlled and fixed
against our experimental groups. First, we can see that Hillis--Steele
\emph{decreases} performance relative to serial when the thread count is limited
to CPU capacities. This is expected: Hillis--Steele is known as a
``work-inefficient'' algorithm, which makes up for its extra work by spreading
it across a hypothetically large number of SIMD-style processing elements. In
this case, since we cap at only 64 threads, the extra work is simply multiplied
across each core, giving $O(n \log n)$ total work and a performance of only
$0.551\times$ the serial baseline at 64 threads.

Analyzing the Blelloch scan, the results stand out as quite interesting. Being
referred to as ``work efficient,'' the Blelloch scan performs $2(n-1)$ total
additions, and, as mentioned earlier, the tree it builds has a depth of
$\log_2(n)$. In our graph, performance increases as the thread count grows and
levels off, reaching up to a $1.75\times$ speedup. The apparent logarithmic
trend might be a coincidence of simply scaling the thread count by powers of two,
but I found it to be an interesting insight into the results as the depth of the work tree is $\log_2(n)$, potentially highlighting a bottleneck for parallelization in this algorithm. Regardless, I
noticed that these results appear bottlenecked by memory bandwidth, and further
investigation into this algorithm would involve a breakdown of how memory
bandwidth affects timing.

\end{document}
